% VI23_FULL_SOLUTION_REFINEMENT_V1
\subsection*{VI.23.P18: A fiber product can be empty}
\begin{problem}\label{prob:vi23-p18}
Give an affine criterion for $\Spec B\times_{\Spec A}\Spec C$ to be empty.
\end{problem}

\ifdefined\FullProblemDossiers
\paragraph{Why this problem matters.}
The criterion shows that nonempty factors need not have any compatible point over the base.

\paragraph{Useful ideas before starting.}
Translate the requested statement into compatible maps from an arbitrary test scheme, and use uniqueness before choosing coordinates.

\paragraph{Guided hints.}
\begin{enumerate}
\item Use the affine product formula.
\item Recall when a spectrum is empty.
\item Construct a concrete example over $k\times k$.
\end{enumerate}

\begin{solution}
The affine fiber product is
\[
\Spec(B\otimes_AC).
\]
For any unital ring $R$, $\Spec R$ is empty exactly when $R$ is the zero ring.  Indeed every
nonzero unital ring has a maximal ideal, hence a prime ideal.  Therefore
\[
\Spec B\times_{\Spec A}\Spec C=\varnothing
\quad\Longleftrightarrow\quad
B\otimes_AC=0.
\]

This can occur even when both factors are nonempty.  For example let
$A=k\times k$, let $B=k$ via the first projection, and let $C=k$ via the second projection.
Then $B$ and $C$ are nonzero fields, but
\[
B\otimes_A C=0:
\]
the idempotent $(1,0)$ acts as $1$ on $B$ and as $0$ on $C$, so for every pure tensor
$b\otimes c=b(1,0)\otimes c=b\otimes (1,0)c=0$.  Thus compatibility over a disconnected base can
eliminate all product points.
\end{solution}

\paragraph{What happened mathematically.}
A geometric assertion was reduced to a representability statement, so the canonical map was forced by universality.

\paragraph{Extensions.}
Repeat the argument after restricting to affine opens and compare with the tensor-product formula.

\paragraph{Provenance.}
\textbf{Canonical tensor-product diagnostic.}
\else
\paragraph{Provenance.} \textbf{Canonical tensor-product diagnostic.}
\fi
